% !TEX root = ../report.tex

\rpzheading{РЕФЕРАТ}

\providecommand{\rpzextraaccountedpages}{0}

Расчётно-пояснительная записка содержит \number\numexpr\ztotpages+\rpzextraaccountedpages\relax{} страницы, \totalfigures{} рисунков, \totaltables{} таблиц, одно приложение, 19 использованных источников.

Ключевые слова: онлайн-маркетплейс, база данных, реляционная модель, postgresql, redis, кэширование, индекс, триггер, хранимая функция, rest api, нагрузочное тестирование.

Объектом разработки является база данных онлайн-маркетплейса и приложение доступа к ней.

Цель работы --- разработать базу данных и приложение для управления товарами, заказами, пользователями, отзывами и отчётами продаж продавца.

В работе проведён анализ предметной области онлайн-маркетплейсов, сформулированы требования к проектируемой системе, выделены роли пользователей и сущности предметной области. Спроектирована реляционная схема базы данных, ролевая модель доступа, триггеры журнала изменения цен и пересчёта рейтингов, а также хранимая функция расчёта статистики продаж продавца.

В реализации используются PostgreSQL, Go и Redis. СУБД PostgreSQL является основным хранилищем данных. Redis используется как in-memory кэш часто читаемых данных и как хранилище отозванных JWT-токенов. Интерфейс доступа включает REST API и серверный Web-интерфейс.

В исследовательской части измерено влияние индексов PostgreSQL и Redis-кэширования на время доступа к данным. Результаты показывают, что индекс даёт эффект только при соответствии его структуры запросу, а прогретый Redis-кэш снижает время HTTP-ответа за счёт сокращения обращений к PostgreSQL.

Область применения разработки --- учебные и демонстрационные системы онлайн-торговли, в которых требуется хранить каталог товаров, заказы, отзывы и отчётность продавца. Дальнейшее развитие может включать расширение аналитических отчётов и административных сценариев.

\clearpage
